\section{Results}\label{sec:results}
\subsection{Main Results}
\begin{table}[htbp]
\centering\small
\caption{Out-of-Sample Forecast Performance}\label{tab:forecast_summary}
\textit{Panel A: 1-Hour RV}\par\smallskip
\begin{adjustbox}{max width=\textwidth}\begin{tabular}{llrrrr}
\toprule
Asset & Model & $R^2$ & DA & QLIKE & MSE \\
\midrule
\multirow{5}{*}{\textbf{BTC}} & HAR & -0.139 & 0.508 & 1.136 & 2.338 \\
 & EGARCH & 0.050 & 0.634 & 0.793 & 1.915 \\
 & EGARCH+NED & 0.336 & 0.703 & 1.153 & 1.338 \\
 & TS-only & 0.417 & 0.733 & 0.992 & 1.175 \\
 & NED & 0.423 & 0.734 & 0.919$^{***}$\,(-7.4\%) & 1.163$^{***}$\,(-1.0\%) \\
\midrule
\multirow{5}{*}{\textbf{ETH}} & HAR & -0.224 & 0.542 & 1.186 & 2.257 \\
 & EGARCH & -0.032 & 0.611 & 0.800 & 1.849 \\
 & EGARCH+NED & 0.253 & 0.681 & 1.203 & 1.335 \\
 & TS-only & 0.335 & 0.710 & 1.038 & 1.188 \\
 & NED & 0.335 & 0.710 & 1.024$^{***}$\,(-1.3\%) & 1.188\,(-0.015\%) \\
\bottomrule
\end{tabular}
\end{adjustbox}
\par\medskip
\textit{Panel B: 4-Hour RV}\par\smallskip
\begin{adjustbox}{max width=\textwidth}\begin{tabular}{llrrrr}
\toprule
Asset & Model & $R^2$ & DA & QLIKE & MSE \\
\midrule
\multirow{5}{*}{\textbf{BTC}} & HAR & -0.232 & 0.500 & 0.711 & 1.769 \\
 & EGARCH & 0.191 & 0.660 & 0.576 & 1.127 \\
 & EGARCH+NED & 0.334 & 0.695 & 0.723 & 0.930 \\
 & TS-only & 0.419 & 0.717 & 0.638 & 0.810 \\
 & NED & 0.438 & 0.723 & 0.569$^{***}$\,(-10.8\%) & 0.785$^{***}$\,(-3.1\%) \\
\midrule
\multirow{5}{*}{\textbf{ETH}} & HAR & -0.328 & 0.547 & 0.727 & 1.653 \\
 & EGARCH & 0.134 & 0.654 & 0.549 & 1.025 \\
 & EGARCH+NED & 0.257 & 0.680 & 0.720 & 0.878 \\
 & TS-only & 0.350 & 0.708 & 0.623 & 0.770 \\
 & NED & 0.350 & 0.710 & 0.610$^{***}$\,(-2.0\%) & 0.771\,(+0.1\%) \\
\bottomrule
\end{tabular}
\end{adjustbox}
\tabnote{\textit{Notes:} Arithmetic means across six expanding windows
(January 2023--December 2025). $R^2$ and MSE use log variance; QLIKE
uses levels. DA is the test-median classification measure in
Equation~\eqref{eq:da}. Parentheses on the NED rows give percentage
changes relative to TS-only; negative values indicate improvement.
$^*p<0.10$, $^{**}p<0.05$, $^{***}p<0.01$ from the corrected DM test
in Equation~\eqref{eq:dm}, applied to concatenated loss differentials.}
\end{table}

Table~\ref{tab:forecast_summary} reports forecast performance at both
horizons for Bitcoin and Ethereum. We begin with the models that do
not use news. The adapted HAR model has negative mean-window $R^2$
at both horizons for both assets. EGARCH improves on this raw-variance
HAR specification, and the neural TS-only model further lowers
log-MSE. At the BTC four-hour horizon, TS-only attains
$R^2=0.419$, compared with 0.191 for EGARCH. The nonlinear model
therefore provides a demanding baseline against which to assess
the incremental value of a news correction.

Adding NED produces significant BTC gains on both forecast losses.
At four hours, QLIKE falls from 0.63815 to 0.56902, a reduction of
10.83\%, while log-MSE falls from 0.80974 to 0.78453, a reduction
of 3.11\%. At one hour, the respective improvements are 7.38\%
and 0.99\%. All four aggregate BTC comparisons have $p<0.001$
under the corrected DM test. Improvements in $R^2$ are more modest,
from 0.41674 to 0.42278 at one hour and from 0.41932 to 0.43759
at four hours. The stronger QLIKE response indicates that the news
correction affects the calibration of variance levels beyond its
smaller change in average squared log error.

Ethereum provides a less uniform result. QLIKE improves by 1.32\%
at one hour and 2.04\% at four hours, with both aggregate comparisons
significant at the 1\% level. Log-MSE is almost unchanged: it falls
by about 0.015\% at one hour and rises by 0.128\% at four hours,
with neither difference statistically significant. Thus the BTC
finding extends to ETH under QLIKE, but not under both losses.

Comparing NED with EGARCH+NED clarifies the role of the baseline.
The hybrid lowers EGARCH log-MSE in all four asset--horizon settings.
For BTC four-hour forecasts, log-MSE falls from 1.12726 to 0.92985
and $R^2$ rises from 0.19067 to 0.33382. These proportional MSE
gains are larger than those obtained by correcting the stronger
neural baseline. However, EGARCH+NED worsens QLIKE in every setting:
BTC four-hour QLIKE rises from 0.57616 to 0.72304. The forecast
value of the correction is therefore loss-dependent, rather than
an improvement that applies to every measure and every baseline.

Standalone EGARCH remains the lowest-QLIKE model for BTC one hour
and for both ETH horizons. NED has the lowest QLIKE for BTC four
hours and the lowest log-MSE in three of the four settings; TS-only
has slightly lower log-MSE for ETH four hours. These rankings
underscore the importance of evaluating more than one loss function.
A model trained on robust log errors need not minimize variance-ratio
loss, even when it improves on an identical time-series checkpoint.

Per-window results in Tables~\ref{tab:expanding_btc}
and~\ref{tab:expanding_eth} show that BTC gains are positive in all
six windows at both horizons and under both losses. The largest
four-hour QLIKE gain, 24.62\%, occurs in the first half of 2024.
ETH improves in QLIKE over the first four windows, but the gain
reverses in both 2025 windows: QLIKE worsens by 3.34\% and 3.17\%
at one hour, and by 4.73\% and 4.85\% at four hours. The aggregate
ETH finding consequently masks economically relevant variation
across periods. Unequal news coverage may be one explanation,
but these comparisons do not identify the cause of the asset difference.

\subsection{Z-score analysis}
We next ask whether the contribution of news is concentrated in
large return-shock regimes. For every BTC four-hour forecast, define
\begin{equation}
 z_t=\frac{\sum_{j=1}^{16}r_{t+j}}
            {\sqrt{\widehat{\RV}_{t,16}^{\mathrm{TS}}}}.
\end{equation}
This is the forward return standardized by the TS-only forecast
volatility. We pool the test observations and form equal-frequency
deciles of this signed score. The grouping distinguishes unusually
large negative and positive returns from observations near the
baseline's predicted scale.

\begin{table}[htbp]
\centering\small
\caption{Forecast performance by standardized-return decile (4-Hour RV, BTC)}
\label{tab:decile_performance}
\begin{adjustbox}{max width=\textwidth}\begin{tabular}{rrrrrrr}
\toprule
Decile & $\bar z$ & TS-only $R^2$ & NED $R^2$ & $\Delta R^2$ & QLIKE (\%) & MSE (\%) \\
\midrule
1 & -2.28 & -0.235 & -0.077 & 0.158 & -16.8 & -12.8 \\
2 & -0.94 & 0.611 & 0.630 & 0.019 & -13.5 & -4.9 \\
3 & -0.55 & 0.568 & 0.555 & -0.013 & -3.3 & +3.0 \\
4 & -0.29 & 0.492 & 0.470 & -0.022 & -1.7 & +4.3 \\
5 & -0.08 & 0.427 & 0.405 & -0.021 & -1.4 & +3.7 \\
6 & 0.13 & 0.446 & 0.426 & -0.021 & -1.2 & +3.7 \\
7 & 0.35 & 0.485 & 0.462 & -0.023 & -1.1 & +4.4 \\
8 & 0.61 & 0.511 & 0.491 & -0.020 & -3.5 & +4.1 \\
9 & 1.00 & 0.569 & 0.577 & 0.008 & -10.1 & -1.9 \\
10 & 2.37 & -0.384 & -0.227 & 0.156 & -12.9 & -11.3 \\
\bottomrule
\end{tabular}
\end{adjustbox}
\tabnote{\textit{Notes:} Deciles use the signed forward-return score
$z_t$, not standardized realized variance. Each contains 10,457
observations. $R^2$ is computed within each pooled decile on log RV;
$\Delta R^2$ is NED minus TS-only. Percentage changes are relative
to TS-only, with negative values indicating lower loss. The grouping
is an ex-post diagnostic, not a real-time decision rule.}
\end{table}

Table~\ref{tab:decile_performance} shows improvements at both tails.
In the lowest decile, with mean $z=-2.28$, NED reduces QLIKE by
16.77\% and log-MSE by 12.76\%. In the highest decile, with mean
$z=2.37$, the corresponding reductions are 12.91\% and 11.29\%.
The $R^2$ gain is about 0.16 in each tail. These results are
consistent with the correction helping when realized returns are
large relative to the baseline forecast.

The middle deciles qualify this interpretation. QLIKE improves
throughout the distribution, but log-MSE rises in deciles three
through eight. News therefore sharpens some aspects of level
calibration while trading off squared log accuracy in more typical
return states. Since the decile labels use future returns, this
analysis describes where the model helped ex post; it does not
establish that favorable states could have been selected in advance.
